% Options for packages loaded elsewhere
\PassOptionsToPackage{unicode}{hyperref}
\PassOptionsToPackage{hyphens}{url}
\documentclass[
  a4paper]{article}
\usepackage{xcolor}
\usepackage[margin=25mm]{geometry}
\usepackage{amsmath,amssymb}
\setcounter{secnumdepth}{-\maxdimen} % remove section numbering
\usepackage{iftex}
\ifPDFTeX
  \usepackage[T1]{fontenc}
  \usepackage[utf8]{inputenc}
  \usepackage{textcomp} % provide euro and other symbols
\else % if luatex or xetex
  \usepackage{unicode-math} % this also loads fontspec
  \defaultfontfeatures{Scale=MatchLowercase}
  \defaultfontfeatures[\rmfamily]{Ligatures=TeX,Scale=1}
\fi
\usepackage{lmodern}
\ifPDFTeX\else
  % xetex/luatex font selection
\fi
% Use upquote if available, for straight quotes in verbatim environments
\IfFileExists{upquote.sty}{\usepackage{upquote}}{}
\IfFileExists{microtype.sty}{% use microtype if available
  \usepackage[]{microtype}
  \UseMicrotypeSet[protrusion]{basicmath} % disable protrusion for tt fonts
}{}
\makeatletter
\@ifundefined{KOMAClassName}{% if non-KOMA class
  \IfFileExists{parskip.sty}{%
    \usepackage{parskip}
  }{% else
    \setlength{\parindent}{0pt}
    \setlength{\parskip}{6pt plus 2pt minus 1pt}}
}{% if KOMA class
  \KOMAoptions{parskip=half}}
\makeatother
\usepackage{color}
\usepackage{fancyvrb}
\newcommand{\VerbBar}{|}
\newcommand{\VERB}{\Verb[commandchars=\\\{\}]}
\DefineVerbatimEnvironment{Highlighting}{Verbatim}{commandchars=\\\{\}}
% Add ',fontsize=\small' for more characters per line
\newenvironment{Shaded}{}{}
\newcommand{\AlertTok}[1]{\textcolor[rgb]{1.00,0.00,0.00}{\textbf{#1}}}
\newcommand{\AnnotationTok}[1]{\textcolor[rgb]{0.38,0.63,0.69}{\textbf{\textit{#1}}}}
\newcommand{\AttributeTok}[1]{\textcolor[rgb]{0.49,0.56,0.16}{#1}}
\newcommand{\BaseNTok}[1]{\textcolor[rgb]{0.25,0.63,0.44}{#1}}
\newcommand{\BuiltInTok}[1]{\textcolor[rgb]{0.00,0.50,0.00}{#1}}
\newcommand{\CharTok}[1]{\textcolor[rgb]{0.25,0.44,0.63}{#1}}
\newcommand{\CommentTok}[1]{\textcolor[rgb]{0.38,0.63,0.69}{\textit{#1}}}
\newcommand{\CommentVarTok}[1]{\textcolor[rgb]{0.38,0.63,0.69}{\textbf{\textit{#1}}}}
\newcommand{\ConstantTok}[1]{\textcolor[rgb]{0.53,0.00,0.00}{#1}}
\newcommand{\ControlFlowTok}[1]{\textcolor[rgb]{0.00,0.44,0.13}{\textbf{#1}}}
\newcommand{\DataTypeTok}[1]{\textcolor[rgb]{0.56,0.13,0.00}{#1}}
\newcommand{\DecValTok}[1]{\textcolor[rgb]{0.25,0.63,0.44}{#1}}
\newcommand{\DocumentationTok}[1]{\textcolor[rgb]{0.73,0.13,0.13}{\textit{#1}}}
\newcommand{\ErrorTok}[1]{\textcolor[rgb]{1.00,0.00,0.00}{\textbf{#1}}}
\newcommand{\ExtensionTok}[1]{#1}
\newcommand{\FloatTok}[1]{\textcolor[rgb]{0.25,0.63,0.44}{#1}}
\newcommand{\FunctionTok}[1]{\textcolor[rgb]{0.02,0.16,0.49}{#1}}
\newcommand{\ImportTok}[1]{\textcolor[rgb]{0.00,0.50,0.00}{\textbf{#1}}}
\newcommand{\InformationTok}[1]{\textcolor[rgb]{0.38,0.63,0.69}{\textbf{\textit{#1}}}}
\newcommand{\KeywordTok}[1]{\textcolor[rgb]{0.00,0.44,0.13}{\textbf{#1}}}
\newcommand{\NormalTok}[1]{#1}
\newcommand{\OperatorTok}[1]{\textcolor[rgb]{0.40,0.40,0.40}{#1}}
\newcommand{\OtherTok}[1]{\textcolor[rgb]{0.00,0.44,0.13}{#1}}
\newcommand{\PreprocessorTok}[1]{\textcolor[rgb]{0.74,0.48,0.00}{#1}}
\newcommand{\RegionMarkerTok}[1]{#1}
\newcommand{\SpecialCharTok}[1]{\textcolor[rgb]{0.25,0.44,0.63}{#1}}
\newcommand{\SpecialStringTok}[1]{\textcolor[rgb]{0.73,0.40,0.53}{#1}}
\newcommand{\StringTok}[1]{\textcolor[rgb]{0.25,0.44,0.63}{#1}}
\newcommand{\VariableTok}[1]{\textcolor[rgb]{0.10,0.09,0.49}{#1}}
\newcommand{\VerbatimStringTok}[1]{\textcolor[rgb]{0.25,0.44,0.63}{#1}}
\newcommand{\WarningTok}[1]{\textcolor[rgb]{0.38,0.63,0.69}{\textbf{\textit{#1}}}}
\setlength{\emergencystretch}{3em} % prevent overfull lines
\providecommand{\tightlist}{%
  \setlength{\itemsep}{0pt}\setlength{\parskip}{0pt}}
\usepackage{bookmark}
\IfFileExists{xurl.sty}{\usepackage{xurl}}{} % add URL line breaks if available
\urlstyle{same}
\hypersetup{
  hidelinks,
  pdfcreator={LaTeX via pandoc}}

\author{}
\date{}

\begin{document}

\section{Anonymous Equal-Amplitude Composite-Wave Model: Basic Axiom
System v3 Interpretation
Note}\label{anonymous-equal-amplitude-composite-wave-model-basic-axiom-system-v3-interpretation-note}

\textbf{Date:} 2026-07-10\\
\textbf{Author:} Noriaki Kihara\\
\textbf{Status:} Interpretation and working note for Basic Axiom System
v3 \textbf{Version DOI:} 10.5281/zenodo.21315736\\
\textbf{Concept DOI:} 10.5281/zenodo.21315735

This document cites the definition paper ``Anonymous Equal-Amplitude
Composite-Wave Model: Basic Axiom System v3'' and records the present
interpretation, development, and working explanations of that axiom
system.

The canonical axiom definition is placed in the following definition
paper:

\begin{Shaded}
\begin{Highlighting}[]
\NormalTok{anonymous\_equal\_amplitude\_composite\_wave\_model\_basic\_axiom\_system\_v3\_pure\_definition\_en.md}
\end{Highlighting}
\end{Shaded}

The interpretations, explanations, working hypotheses, and readout
examples in this document may be extended or revised.

However, a change in interpretation is not a rejection of the axioms.

If an error, omission, or required change is found in the axioms
themselves, the canonical axiom system must be explicitly revised rather
than being reinterpreted inside this note.

Thus this document is not the axiom system itself. It is the current
interpretation layer for Basic Axiom System v3.

In v3, anonymity is extended beyond basic components, physical
quantities, and observables to theory names, formulation names, and
readout-logic names. In other words, the names gravity, Coulomb force,
mass, charge, energy, two-body, three-body, and similar labels must not
be used as reasons to change equations, readout rules, or projection
rules at the first-principle layer. Physical names are labels assigned
after readout results have been obtained from the same axioms.

\begin{center}\rule{0.5\linewidth}{0.5pt}\end{center}

\section{1. First Principles}\label{first-principles}

\subsection{Axiom 0: Anonymity}\label{axiom-0-anonymity}

Classification: Axiom

Basic components are not assigned individual names, privileged axes,
privileged signs, or privileged types. The index \texttt{n} is a
convenient label and is not an intrinsic name of the component.

Forbidden:

\begin{Shaded}
\begin{Highlighting}[]
\NormalTok{making only one specific component the negative{-}sign axis}
\NormalTok{making only one specific component the time axis}
\NormalTok{making only one specific component real{-}valued}
\NormalTok{making only one specific component complex{-}valued}
\NormalTok{giving an intrinsic name to only one specific component}
\end{Highlighting}
\end{Shaded}

Therefore,

\begin{Shaded}
\begin{Highlighting}[]
\NormalTok{x\_1\^{}2+x\_2\^{}2{-}x\_3\^{}2=0}
\end{Highlighting}
\end{Shaded}

is not adopted at the first-principle layer, because \texttt{x\_3} alone
becomes the negative-sign axis.

\subsection{Axiom 0+: Formulation
Anonymity}\label{axiom-0-formulation-anonymity}

Classification: Axiom

At the first-principle layer, theory names, formulation names, and
readout-logic names are not assigned intrinsic names.

Equations, projections, or readout rules at the first-principle layer
must not be changed on the basis of physical names, object names,
interaction names, or existing theory classifications.

Forbidden:

\begin{Shaded}
\begin{Highlighting}[]
\NormalTok{adopting a circular{-}motion equation because it is called gravity}
\NormalTok{adopting a hyperbolic equation because it is called Coulomb force}
\NormalTok{erasing signs because it is called mass}
\NormalTok{preserving signs because it is called charge}
\NormalTok{moving a component to a time axis because it is called energy}
\NormalTok{using a different closure equation because the target is a two{-}body system}
\NormalTok{using a different closure equation because the target is a three{-}body system}
\NormalTok{using a special readout rule because the target is an ABC system}
\NormalTok{choosing a right{-}hand{-}side representation to obtain a desired physical name}
\end{Highlighting}
\end{Shaded}

At the first-principle layer, only the same closure condition,

\begin{Shaded}
\begin{Highlighting}[]
\NormalTok{\textbackslash{}sum\_n x\_n\^{}2=0}
\end{Highlighting}
\end{Shaded}

and readout operations, symmetry breaking, gauge stability, sign
preservation, sign erasure, phase-branch selection, and closure target
sets defined in advance for that same closure condition are allowed.

The following order is forbidden:

\begin{Shaded}
\begin{Highlighting}[]
\NormalTok{call it gravity}
\NormalTok{therefore adopt an R{-}type readout}
\NormalTok{therefore erase signs}
\end{Highlighting}
\end{Shaded}

Likewise, the following order is forbidden:

\begin{Shaded}
\begin{Highlighting}[]
\NormalTok{call it Coulomb force}
\NormalTok{therefore adopt a Q{-}type readout}
\NormalTok{therefore preserve signs}
\end{Highlighting}
\end{Shaded}

The correct order is:

\begin{Shaded}
\begin{Highlighting}[]
\NormalTok{place the same all{-}positive{-}sign zero closure}
\NormalTok{define the readout operation first}
\NormalTok{define symmetry breaking first}
\NormalTok{define sign preservation or sign erasure first}
\NormalTok{confirm gauge stability}
\NormalTok{assign working names after readout}
\end{Highlighting}
\end{Shaded}

Here \texttt{R}, \texttt{Q}, and \texttt{E} are not first-principle
names. They are working names assigned to readout results obtained from
the same closure condition.

In particular, an already defined \texttt{R\_read} must not be changed
into another definition because one wishes to read gravity.

What is allowed is to apply the same \texttt{R\_read} to a single closed
wave, a pair composite closed wave, or a whole-system closed wave, and
to read differences in the target set as:

\begin{Shaded}
\begin{Highlighting}[]
\NormalTok{R\_A}
\NormalTok{R\_\{AB\}}
\NormalTok{R\_\{ABC\}}
\end{Highlighting}
\end{Shaded}

Likewise, the distinction between a two-body and a three-body system is
treated as a difference in closure target set, not as a
physical-name-based change in equation.

For example, the following are allowed because the target sets differ:

\begin{Shaded}
\begin{Highlighting}[]
\NormalTok{Q(A)=0}
\end{Highlighting}
\end{Shaded}

\begin{Shaded}
\begin{Highlighting}[]
\NormalTok{Q(A+B)=0}
\end{Highlighting}
\end{Shaded}

\begin{Shaded}
\begin{Highlighting}[]
\NormalTok{Q(A+B+C)=0}
\end{Highlighting}
\end{Shaded}

However, each \texttt{Q} must have the same form:

\begin{Shaded}
\begin{Highlighting}[]
\NormalTok{Q(P)=\textbackslash{}sum\_n p\_n\^{}2}
\end{Highlighting}
\end{Shaded}

Different readouts may be obtained from different target sets. Changing
the definition of \texttt{Q} itself to obtain a desired physical name is
not allowed.

Therefore, formulation anonymity extends Axiom 0, component anonymity,
to the level of theoretical construction.

This axiom requires:

\begin{Shaded}
\begin{Highlighting}[]
\NormalTok{closure equation before physical name}
\NormalTok{readout operation before classification name}
\NormalTok{symmetry breaking before theory name}
\NormalTok{gauge stability before interpretation name}
\end{Highlighting}
\end{Shaded}

Physical names, classification names, theory names, and interpretation
names are not first principles. They are names assigned after readout.

\subsection{Axiom 1: All-Positive-Sign Zero
Closure}\label{axiom-1-all-positive-sign-zero-closure}

Classification: Axiom

\begin{Shaded}
\begin{Highlighting}[]
\NormalTok{\textbackslash{}sum\_\{n=1\}\^{}\{N\}x\_n\^{}2=0}
\end{Highlighting}
\end{Shaded}

Expanded form:

\begin{Shaded}
\begin{Highlighting}[]
\NormalTok{x\_1\^{}2+x\_2\^{}2+\textbackslash{}cdots+x\_N\^{}2=0}
\end{Highlighting}
\end{Shaded}

All coefficient signs are the same positive sign. This is not the
conjugate norm.

\begin{Shaded}
\begin{Highlighting}[]
\NormalTok{\textbackslash{}sum\_n |x\_n|\^{}2=0}
\end{Highlighting}
\end{Shaded}

is not the intended closure. The closure uses \texttt{x\_n\^{}2}, not
\texttt{x\_n\textbackslash{}bar\{x\}\_n}.

Axiom 1 is the closure condition for a state in which the left-hand side
of the closure sum has become stationary or stable.

It does not deny the existence of a quasi-stationary or metastable state
between an external influence, interaction, observation, or readout
re-embedding that changes the closure sum and the later return to
stationarity or stability.

Therefore, in a quasi-stationary state, the following temporary nonzero
closure residual is allowed:

\begin{Shaded}
\begin{Highlighting}[]
\NormalTok{\textbackslash{}sum\_n x\_n\^{}2\textbackslash{}neq0}
\end{Highlighting}
\end{Shaded}

However, in a state treated as stationary closure or stable closure,

\begin{Shaded}
\begin{Highlighting}[]
\NormalTok{\textbackslash{}sum\_n x\_n\^{}2=0}
\end{Highlighting}
\end{Shaded}

holds.

Thus, later papers may treat closure recovery processes, metastable
transitions, or closure residuals immediately after observation without
rejecting Axiom 1.

Axiom 1 itself defines the closure condition after stationarization, not
the quasi-stationary process.

\subsection{Axiom 2: Nontrivial
Existence}\label{axiom-2-nontrivial-existence}

Classification: Axiom

\begin{Shaded}
\begin{Highlighting}[]
\NormalTok{\textbackslash{}exists n,\textbackslash{}quad x\_n\textbackslash{}neq0}
\end{Highlighting}
\end{Shaded}

The closed system is not restricted to the all-zero solution.

\begin{center}\rule{0.5\linewidth}{0.5pt}\end{center}

\section{2. Consequences from the First
Principles}\label{consequences-from-the-first-principles}

\subsection{Consequence 1: Insufficiency of Real Positive-Definite
Components}\label{consequence-1-insufficiency-of-real-positive-definite-components}

Classification: Derived consequence

\begin{Shaded}
\begin{Highlighting}[]
\NormalTok{x\_n\textbackslash{}in\textbackslash{}mathbb\{R\}}
\NormalTok{\textbackslash{}Rightarrow}
\NormalTok{x\_n\^{}2\textbackslash{}ge0}
\end{Highlighting}
\end{Shaded}

Therefore,

\begin{Shaded}
\begin{Highlighting}[]
\NormalTok{\textbackslash{}sum\_n x\_n\^{}2=0}
\NormalTok{\textbackslash{}Rightarrow}
\NormalTok{x\_n=0\textbackslash{}quad\textbackslash{}forall n}
\end{Highlighting}
\end{Shaded}

To satisfy Axioms 0, 1, and 2 simultaneously, real positive-definite
components alone are insufficient.

\subsection{Consequence 2: Internal Generation of Sign
Reversal}\label{consequence-2-internal-generation-of-sign-reversal}

Classification: Derived consequence

A negative sign as an external coefficient violates Axiom 0. The sign
reversal required for closure must arise from the square of a component.

Minimal two-component solution:

\begin{Shaded}
\begin{Highlighting}[]
\NormalTok{x\_1=A,\textbackslash{}qquad x\_2=iA}
\end{Highlighting}
\end{Shaded}

\begin{Shaded}
\begin{Highlighting}[]
\NormalTok{x\_1\^{}2+x\_2\^{}2=A\^{}2+(iA)\^{}2=A\^{}2{-}A\^{}2=0}
\end{Highlighting}
\end{Shaded}

\begin{Shaded}
\begin{Highlighting}[]
\NormalTok{i\^{}2={-}1}
\end{Highlighting}
\end{Shaded}

The negative sign is not a privileged coefficient. It is a squared
phase.

\subsection{Consequence 3: Necessity of Phase
Algebra}\label{consequence-3-necessity-of-phase-algebra}

Classification: Derived consequence

Nontrivial closure requires a common phase algebra that internally
generates sign reversal by squaring.

\begin{Shaded}
\begin{Highlighting}[]
\NormalTok{x\_n=r\_n e\^{}\{i\textbackslash{}phi\_n\}}
\end{Highlighting}
\end{Shaded}

\begin{Shaded}
\begin{Highlighting}[]
\NormalTok{x\_n\^{}2=r\_n\^{}2e\^{}\{i2\textbackslash{}phi\_n\}}
\end{Highlighting}
\end{Shaded}

Closure condition:

\begin{Shaded}
\begin{Highlighting}[]
\NormalTok{\textbackslash{}sum\_\{n=1\}\^{}\{N\}r\_n\^{}2e\^{}\{i2\textbackslash{}phi\_n\}=0}
\end{Highlighting}
\end{Shaded}

Under the equal-amplitude condition \texttt{r\_n=r},

\begin{Shaded}
\begin{Highlighting}[]
\NormalTok{\textbackslash{}sum\_\{n=1\}\^{}\{N\}e\^{}\{i2\textbackslash{}phi\_n\}=0}
\end{Highlighting}
\end{Shaded}

The complex phase type is the minimal implementation of this
requirement.

\subsection{Consequence 4: Type
Anonymity}\label{consequence-4-type-anonymity}

Classification: Derived consequence

Making only one component complex-valued violates Axiom 0. Therefore,

\begin{Shaded}
\begin{Highlighting}[]
\NormalTok{x\_n\textbackslash{}in\textbackslash{}mathbb\{C\}\textbackslash{}quad\textbackslash{}forall n}
\end{Highlighting}
\end{Shaded}

or all components must be given the same phase-algebra type.

\subsection{Consequence 5: Ninety-Degree Phase
Difference}\label{consequence-5-ninety-degree-phase-difference}

Classification: Derived consequence

\begin{Shaded}
\begin{Highlighting}[]
\NormalTok{A\^{}2+(iA)\^{}2=0}
\end{Highlighting}
\end{Shaded}

A ninety-degree phase difference is the minimal two-component structure
of all-positive-sign zero closure.

\begin{center}\rule{0.5\linewidth}{0.5pt}\end{center}

\section{3. Anonymous Equal-Amplitude Composite
Waves}\label{anonymous-equal-amplitude-composite-waves}

\subsection{Axiom 3: Anonymous Equal-Amplitude Composite
Wave}\label{axiom-3-anonymous-equal-amplitude-composite-wave}

Classification: Axiom

\begin{Shaded}
\begin{Highlighting}[]
\NormalTok{\textbackslash{}Psi\_N=\textbackslash{}frac\{A\_0\}\{N\}\textbackslash{}sum\_\{k=1\}\^{}\{N\}e\^{}\{i\textbackslash{}theta\_k\}}
\end{Highlighting}
\end{Shaded}

\begin{Shaded}
\begin{Highlighting}[]
\NormalTok{|\textbackslash{}psi\_k|=\textbackslash{}frac\{A\_0\}\{N\}}
\end{Highlighting}
\end{Shaded}

Components are not individually identified by amplitude. Since anonymous
basic components are not distinguished by amplitude differences, all
components are taken to have equal amplitude in the basic implementation
of the composite wave.

Normalization condition:

When all phases are aligned,

\begin{Shaded}
\begin{Highlighting}[]
\NormalTok{\textbackslash{}max|\textbackslash{}Psi\_N|=A\_0}
\end{Highlighting}
\end{Shaded}

Thus the normalization is

\begin{Shaded}
\begin{Highlighting}[]
\NormalTok{\textbackslash{}Psi\_N=\textbackslash{}frac\{A\_0\}\{N\}\textbackslash{}sum\_\{k=1\}\^{}\{N\}e\^{}\{i\textbackslash{}theta\_k\}}
\end{Highlighting}
\end{Shaded}

What is preserved is the maximum composite amplitude readable from
outside.

\subsection{Axiom 4: Unreadability of Individual Absolute
Phase}\label{axiom-4-unreadability-of-individual-absolute-phase}

Classification: Axiom

\begin{Shaded}
\begin{Highlighting}[]
\NormalTok{\textbackslash{}theta\_k\textbackslash{}rightarrow\textbackslash{}theta\_k+\textbackslash{}alpha}
\end{Highlighting}
\end{Shaded}

does not change the internal structure. Absolute phase is not an
intrinsic observable.

\subsection{Axiom 5: Phase-Difference
Readout}\label{axiom-5-phase-difference-readout}

Classification: Axiom

\begin{Shaded}
\begin{Highlighting}[]
\NormalTok{\textbackslash{}Delta\_\{ij\}=\textbackslash{}theta\_i{-}\textbackslash{}theta\_j}
\end{Highlighting}
\end{Shaded}

The number of readable pair channels is

\begin{Shaded}
\begin{Highlighting}[]
\NormalTok{\textbackslash{}binom\{N\}\{2\}=\textbackslash{}frac\{N(N{-}1)\}\{2\}}
\end{Highlighting}
\end{Shaded}

This is not the number of degrees of freedom, but the number of
phase-difference relations.

\subsection{Axiom 6: System
Identification}\label{axiom-6-system-identification}

Classification: Axiom

Two systems with the same

\begin{Shaded}
\begin{Highlighting}[]
\NormalTok{(N,A\_0,\textbackslash{}\{\textbackslash{}Delta\_\{ij\}\textbackslash{}\})}
\end{Highlighting}
\end{Shaded}

are identified as the same system.

\subsection{Axiom 7: Composite Phase
Projection}\label{axiom-7-composite-phase-projection}

Classification: Axiom

\begin{Shaded}
\begin{Highlighting}[]
\NormalTok{S\_N=\textbackslash{}sum\_k e\^{}\{i\textbackslash{}theta\_k\}}
\end{Highlighting}
\end{Shaded}

\begin{Shaded}
\begin{Highlighting}[]
\NormalTok{\textbackslash{}Psi\_N=\textbackslash{}frac\{A\_0\}\{N\}S\_N}
\end{Highlighting}
\end{Shaded}

\begin{Shaded}
\begin{Highlighting}[]
\NormalTok{\textbackslash{}Theta\_N=\textbackslash{}arg\textbackslash{}Psi\_N}
\end{Highlighting}
\end{Shaded}

\texttt{Θ\_N} is an external projected phase defined when
\texttt{Ψ\_N\ !=\ 0}. Even when \texttt{Ψ\_N=0}, the internal
phase-difference network exists.

\subsection{Axiom 8: Reflection Symmetry Is Not
Required}\label{axiom-8-reflection-symmetry-is-not-required}

Classification: Axiom

A composite wave is not required to be reflection symmetric.

\subsection{Consequence 6: Reflection Symmetry Is a Special
Condition}\label{consequence-6-reflection-symmetry-is-a-special-condition}

Classification: Derived consequence from Axiom 8

Reflection symmetry is a special condition.

\begin{Shaded}
\begin{Highlighting}[]
\NormalTok{\textbackslash{}Psi({-}x)=\textbackslash{}Psi(x)}
\end{Highlighting}
\end{Shaded}

In general, reflection-asymmetric configurations are allowed.

\begin{Shaded}
\begin{Highlighting}[]
\NormalTok{\textbackslash{}Psi({-}x)\textbackslash{}neq\textbackslash{}Psi(x)}
\end{Highlighting}
\end{Shaded}

A reflection-symmetric wave is a special configuration corresponding to
equal weighting of forward and reverse phase progression.

\begin{center}\rule{0.5\linewidth}{0.5pt}\end{center}

\section{4. External Readout}\label{external-readout}

\subsection{Axiom 9: External-Axis
Response}\label{axiom-9-external-axis-response}

Classification: Candidate axiom

Let the external axis be

\begin{Shaded}
\begin{Highlighting}[]
\NormalTok{M\_\textbackslash{}alpha=e\^{}\{i\textbackslash{}alpha\}}
\end{Highlighting}
\end{Shaded}

Then

\begin{Shaded}
\begin{Highlighting}[]
\NormalTok{Y(\textbackslash{}alpha)=\textbackslash{}frac\{A\_0\}\{N\}\textbackslash{}sum\_\{k=1\}\^{}\{N\}e\^{}\{i(\textbackslash{}theta\_k{-}\textbackslash{}alpha)\}}
\NormalTok{=e\^{}\{{-}i\textbackslash{}alpha\}\textbackslash{}Psi\_N}
\end{Highlighting}
\end{Shaded}

Continuous response:

\begin{Shaded}
\begin{Highlighting}[]
\NormalTok{S\_N(\textbackslash{}alpha)=\textbackslash{}operatorname\{Re\}\textbackslash{}left[\textbackslash{}frac\{1\}\{N\}\textbackslash{}sum\_\{k=1\}\^{}\{N\}e\^{}\{i(\textbackslash{}theta\_k{-}\textbackslash{}alpha)\}\textbackslash{}right]}
\end{Highlighting}
\end{Shaded}

Binary response:

\begin{Shaded}
\begin{Highlighting}[]
\NormalTok{S\_N(\textbackslash{}alpha)=\textbackslash{}operatorname\{sgn\}\textbackslash{}operatorname\{Re\}\textbackslash{}left[\textbackslash{}frac\{1\}\{N\}\textbackslash{}sum\_\{k=1\}\^{}\{N\}e\^{}\{i(\textbackslash{}theta\_k{-}\textbackslash{}alpha)\}\textbackslash{}right]}
\end{Highlighting}
\end{Shaded}

Spin-like information is defined by correlation readout with an external
axis.

\subsection{Axiom 10: External Readout
Phase}\label{axiom-10-external-readout-phase}

Classification: Candidate axiom

\begin{Shaded}
\begin{Highlighting}[]
\NormalTok{\textbackslash{}Delta\textbackslash{}phi\_\textbackslash{}mu=\textbackslash{}phi\_\textbackslash{}mu\^{}\{(\textbackslash{}Psi)\}{-}\textbackslash{}phi\_\textbackslash{}mu\^{}\{(M)\}}
\end{Highlighting}
\end{Shaded}

\begin{Shaded}
\begin{Highlighting}[]
\NormalTok{\textbackslash{}mu=x,y,z,t}
\end{Highlighting}
\end{Shaded}

Position-like, time-like, and spin-like readouts are relative phases
with respect to an external reference.

\subsection{Axiom 11: Position-Phase
Readout}\label{axiom-11-position-phase-readout}

Classification: Candidate axiom

\begin{Shaded}
\begin{Highlighting}[]
\NormalTok{x,y,z}
\end{Highlighting}
\end{Shaded}

are not intrinsic labels of the composite wave. They are reconstructed
from relative phase with an external spatial readout wave.

The position notation used here is a convenient reinterpretation and
does not prove connection to physical position quantities.

\subsection{Axiom 12: Time-Phase
Readout}\label{axiom-12-time-phase-readout}

Classification: Candidate axiom

\begin{Shaded}
\begin{Highlighting}[]
\NormalTok{\textbackslash{}phi\_t\textbackslash{}equiv\textbackslash{}omega t\_\{\textbackslash{}mathrm\{read\}\}\textbackslash{}pmod\{2\textbackslash{}pi\}}
\end{Highlighting}
\end{Shaded}

Time is a relative-phase readout with respect to an external frequency
reference. The parameter \texttt{t\_read} here is not an a priori
absolute time, but a convenient parameter reconstructed from time-phase
readout.

\subsection{Axiom 13: Covering Phase}\label{axiom-13-covering-phase}

Classification: Candidate axiom

When the readout system preserves winding number, it distinguishes:

\begin{Shaded}
\begin{Highlighting}[]
\NormalTok{2π{-}period readout}
\NormalTok{4π{-}period readout}
\NormalTok{winding{-}difference readout}
\end{Highlighting}
\end{Shaded}

\begin{center}\rule{0.5\linewidth}{0.5pt}\end{center}

\section{5. Observed Phase
Discreteness}\label{observed-phase-discreteness}

\subsection{Axiom 14: Observed Phase
Discreteness}\label{axiom-14-observed-phase-discreteness}

Classification: Candidate axiom

A phase fixed as an observed value is discretized by a loop condition or
winding number.

\begin{Shaded}
\begin{Highlighting}[]
\NormalTok{\textbackslash{}sum\_\{\textbackslash{}mathrm\{cycle\}\}\textbackslash{}Delta\textbackslash{}phi=2\textbackslash{}pi m,\textbackslash{}qquad m\textbackslash{}in\textbackslash{}mathbb\{Z\}}
\end{Highlighting}
\end{Shaded}

\begin{Shaded}
\begin{Highlighting}[]
\NormalTok{W=\textbackslash{}frac\{1\}\{2\textbackslash{}pi\}\textbackslash{}oint d\textbackslash{}phi,\textbackslash{}qquad W\textbackslash{}in\textbackslash{}mathbb\{Z\}}
\end{Highlighting}
\end{Shaded}

Readout-cell form:

\begin{Shaded}
\begin{Highlighting}[]
\NormalTok{\textbackslash{}phi\_\{\textbackslash{}mathrm\{obs\}\}=m\textbackslash{}epsilon\_\textbackslash{}phi,\textbackslash{}qquad m\textbackslash{}in\textbackslash{}mathbb\{Z\}}
\end{Highlighting}
\end{Shaded}

Candidate phase-area quantity:

\begin{Shaded}
\begin{Highlighting}[]
\NormalTok{Q\_\textbackslash{}phi=\textbackslash{}sum\_\{i\textless{}j\}\textbackslash{}sin\^{}2\textbackslash{}frac\{\textbackslash{}Delta\_\{ij\}\}\{2\}}
\end{Highlighting}
\end{Shaded}

\subsection{Readout Example: Bounded
Closure}\label{readout-example-bounded-closure}

Classification: Implementation example

In a closed readout system,

\begin{Shaded}
\begin{Highlighting}[]
\NormalTok{kL=2\textbackslash{}pi m}
\end{Highlighting}
\end{Shaded}

\begin{Shaded}
\begin{Highlighting}[]
\NormalTok{k=\textbackslash{}frac\{2\textbackslash{}pi m\}\{L\}}
\end{Highlighting}
\end{Shaded}

The observable phase change is integerized as a winding number. The
first-principle closure form is Axiom 1.

\begin{center}\rule{0.5\linewidth}{0.5pt}\end{center}

\section{6. Waveform, Localization, and
Observation}\label{waveform-localization-and-observation}

\subsection{Working Axiom A: State
Waveform}\label{working-axiom-a-state-waveform}

Classification: Working hypothesis

\begin{Shaded}
\begin{Highlighting}[]
\NormalTok{A(\textbackslash{}theta)=\textbackslash{}sum\_n A\_n\textbackslash{}cos(n\textbackslash{}theta+\textbackslash{}phi\_n)}
\end{Highlighting}
\end{Shaded}

\begin{Shaded}
\begin{Highlighting}[]
\NormalTok{Z(\textbackslash{}theta)=\textbackslash{}sum\_n r\_n e\^{}\{i(n\textbackslash{}theta+\textbackslash{}phi\_n)\}}
\end{Highlighting}
\end{Shaded}

A state is a hierarchical waveform.

\subsection{Definition A1: Normalized Equal-Amplitude Odd-Harmonic
Localized
Wave}\label{definition-a1-normalized-equal-amplitude-odd-harmonic-localized-wave}

Classification: Definition

Let the highest odd-harmonic order be

\begin{Shaded}
\begin{Highlighting}[]
\NormalTok{N\_h=2K{-}1}
\end{Highlighting}
\end{Shaded}

Define the normalized equal-amplitude odd-harmonic localized wave by

\begin{Shaded}
\begin{Highlighting}[]
\NormalTok{S\_\{N\_h\}(u)=\textbackslash{}frac\{1\}\{K\}\textbackslash{}sum\_\{m=0\}\^{}\{K{-}1\}\textbackslash{}cos((2m+1)u)}
\end{Highlighting}
\end{Shaded}

In complex notation,

\begin{Shaded}
\begin{Highlighting}[]
\NormalTok{Z\_\{N\_h\}(u)=\textbackslash{}frac\{1\}\{K\}\textbackslash{}sum\_\{m=0\}\^{}\{K{-}1\}e\^{}\{i(2m+1)u\}}
\end{Highlighting}
\end{Shaded}

A localized wave with representative amplitude \texttt{A} is

\begin{Shaded}
\begin{Highlighting}[]
\NormalTok{\textbackslash{}Psi\_\{N\_h\}(u)=A S\_\{N\_h\}(u)}
\end{Highlighting}
\end{Shaded}

or

\begin{Shaded}
\begin{Highlighting}[]
\NormalTok{\textbackslash{}Psi\_\{N\_h\}(u)=A Z\_\{N\_h\}(u)}
\end{Highlighting}
\end{Shaded}

Each odd-harmonic component has amplitude

\begin{Shaded}
\begin{Highlighting}[]
\NormalTok{\textbackslash{}frac\{A\}\{K\}}
\end{Highlighting}
\end{Shaded}

and is normalized by the inverse of the number \texttt{K} of components.
This normalization keeps the representative amplitude at in-phase
alignment equal to \texttt{A}, even as \texttt{N\_h} increases.

\subsection{Theorem A1: Preservation of Representative
Amplitude}\label{theorem-a1-preservation-of-representative-amplitude}

Classification: Derived theorem

Consider

\begin{Shaded}
\begin{Highlighting}[]
\NormalTok{S\_\{N\_h\}(u)=\textbackslash{}frac\{1\}\{K\}\textbackslash{}sum\_\{m=0\}\^{}\{K{-}1\}\textbackslash{}cos((2m+1)u)}
\end{Highlighting}
\end{Shaded}

At the central phase \texttt{u=0},

\begin{Shaded}
\begin{Highlighting}[]
\NormalTok{\textbackslash{}cos((2m+1)0)=1}
\end{Highlighting}
\end{Shaded}

and hence

\begin{Shaded}
\begin{Highlighting}[]
\NormalTok{S\_\{N\_h\}(0)=\textbackslash{}frac\{1\}\{K\}\textbackslash{}sum\_\{m=0\}\^{}\{K{-}1\}1=1}
\end{Highlighting}
\end{Shaded}

For

\begin{Shaded}
\begin{Highlighting}[]
\NormalTok{\textbackslash{}Psi\_\{N\_h\}(u)=A S\_\{N\_h\}(u)}
\end{Highlighting}
\end{Shaded}

we obtain

\begin{Shaded}
\begin{Highlighting}[]
\NormalTok{\textbackslash{}Psi\_\{N\_h\}(0)=A}
\end{Highlighting}
\end{Shaded}

Therefore, even as the highest odd-harmonic order \texttt{N\_h}
increases, the normalization factor \texttt{1/K} preserves the composite
peak representative amplitude as \texttt{A}.

\subsection{Theorem A2: Highest Odd-Harmonic Order and Localization
Width}\label{theorem-a2-highest-odd-harmonic-order-and-localization-width}

Classification: Derived theorem

Let the base wavelength of the parent system be

\begin{Shaded}
\begin{Highlighting}[]
\NormalTok{\textbackslash{}lambda\_0}
\end{Highlighting}
\end{Shaded}

Here \texttt{N\_h} is not the number of odd-harmonic components, but the
highest odd-harmonic multiplier relative to the parent base wave. The
number of odd-harmonic components is

\begin{Shaded}
\begin{Highlighting}[]
\NormalTok{K=\textbackslash{}frac\{N\_h+1\}\{2\}}
\end{Highlighting}
\end{Shaded}

The wavelength corresponding to the highest harmonic is conceptually

\begin{Shaded}
\begin{Highlighting}[]
\NormalTok{\textbackslash{}lambda\_\{N\_h\}=\textbackslash{}frac\{\textbackslash{}lambda\_0\}\{N\_h\}}
\end{Highlighting}
\end{Shaded}

Thus the localization width obtained by normalized equal-amplitude
odd-harmonic synthesis is estimated as

\begin{Shaded}
\begin{Highlighting}[]
\NormalTok{\textbackslash{}Delta x\_\{N\_h\}\textbackslash{}sim\textbackslash{}frac\{\textbackslash{}lambda\_0\}\{N\_h\}}
\end{Highlighting}
\end{Shaded}

or, using a half-wavelength criterion,

\begin{Shaded}
\begin{Highlighting}[]
\NormalTok{\textbackslash{}Delta x\_\{N\_h\}\textbackslash{}sim\textbackslash{}frac\{\textbackslash{}lambda\_0\}\{2N\_h\}}
\end{Highlighting}
\end{Shaded}

The exact coefficient depends on whether one uses main-lobe width,
half-maximum width, or first-zero width. This document uses

\begin{Shaded}
\begin{Highlighting}[]
\NormalTok{\textbackslash{}Delta x\_\{N\_h\}\textbackslash{}propto\textbackslash{}frac\{\textbackslash{}lambda\_0\}\{N\_h\}}
\end{Highlighting}
\end{Shaded}

as a design estimate.

Therefore, the highest odd-harmonic order required to obtain a local
width \texttt{Δx} from a base wavelength \texttt{λ\_0} is approximately

\begin{Shaded}
\begin{Highlighting}[]
\NormalTok{N\_\{h,\textbackslash{}rm req\}\textbackslash{}sim\textbackslash{}frac\{\textbackslash{}lambda\_0\}\{\textbackslash{}Delta x\}}
\end{Highlighting}
\end{Shaded}

In implementation, an odd order at least as large as
\texttt{N\_\{h,\textbackslash{}rm\ req\}} is used.

For the temporal direction, if the parent base period is \texttt{T\_0}
and the temporal localization width is \texttt{Δt}, then

\begin{Shaded}
\begin{Highlighting}[]
\NormalTok{N\_\{h,\textbackslash{}rm req\}\^{}\{(t)\}\textbackslash{}sim\textbackslash{}frac\{T\_0\}\{\textbackslash{}Delta t\}}
\end{Highlighting}
\end{Shaded}

\subsection{Consequence A1: Common Construction of Spatial and Temporal
Localization}\label{consequence-a1-common-construction-of-spatial-and-temporal-localization}

Classification: Derived consequence

The normalized equal-amplitude odd-harmonic localized wave can be
applied not only to spatial phase but also to temporal phase.

Let spatial phase be \texttt{χ} and temporal phase be \texttt{τ}. A
local wave localized in both spatial and temporal directions may be
written as

\begin{Shaded}
\begin{Highlighting}[]
\NormalTok{\textbackslash{}Psi(\textbackslash{}chi,\textbackslash{}tau)=A S\_\{N\_\{h,\textbackslash{}chi\}\}(\textbackslash{}chi{-}\textbackslash{}chi\_0)S\_\{N\_\{h,\textbackslash{}tau\}\}(\textbackslash{}tau{-}\textbackslash{}tau\_0)e\^{}\{i\textbackslash{}phi\}}
\end{Highlighting}
\end{Shaded}

where \texttt{χ\_0} is position phase, \texttt{τ\_0} is time phase,
\texttt{N\_\{h,χ\}} is the highest odd-harmonic order in the spatial
direction, and \texttt{N\_\{h,τ\}} is the highest odd-harmonic order in
the temporal direction.

This form represents spatial and temporal localization by the
odd-harmonic structure of the waveform itself, without introducing an
externally invisible temporal window function.

\subsection{Note A1: Quantity
Preserved}\label{note-a1-quantity-preserved}

Classification: Note

In this axiom system, what is preserved by the normalized
equal-amplitude odd-harmonic localized wave is the representative
amplitude \texttt{A} and the square register governed by Axiom 1,
all-positive-sign zero closure.

Normalizing each odd-harmonic component to amplitude

\begin{Shaded}
\begin{Highlighting}[]
\NormalTok{\textbackslash{}frac\{A\}\{K\}}
\end{Highlighting}
\end{Shaded}

keeps the in-phase composite representative amplitude at

\begin{Shaded}
\begin{Highlighting}[]
\NormalTok{A}
\end{Highlighting}
\end{Shaded}

Thus, increasing \texttt{N\_h} does not make the in-phase representative
amplitude diverge.

This document places the basis of conservation in

\begin{Shaded}
\begin{Highlighting}[]
\NormalTok{\textbackslash{}sum\_n x\_n\^{}2=0}
\end{Highlighting}
\end{Shaded}

It is not a conjugate-norm conservation law

\begin{Shaded}
\begin{Highlighting}[]
\NormalTok{\textbackslash{}sum\_n |x\_n|\^{}2}
\end{Highlighting}
\end{Shaded}

and it is not a conservation law with externally imposed weights on
higher harmonic orders.

Localization by higher harmonics is treated as refinement of phase
structure while preserving the representative amplitude \texttt{A}.
Increasing the highest odd-harmonic order is therefore not an operation
that externally increases the conserved quantity. It is an operation
that refines phase structure with the same representative amplitude
inside the closure condition of Axiom 1.

\subsection{Working Axiom B:
Localization}\label{working-axiom-b-localization}

Classification: Working hypothesis

Localization is represented as phase alignment of higher odd-harmonic
components with respect to a low-order base wave. This document does not
introduce localization as an external window function.

Localization is constructed by

\begin{Shaded}
\begin{Highlighting}[]
\NormalTok{S\_\{N\_h\}(u)=\textbackslash{}frac\{1\}\{K\}\textbackslash{}sum\_\{m=0\}\^{}\{K{-}1\}\textbackslash{}cos((2m+1)u)}
\end{Highlighting}
\end{Shaded}

or by its complex representation

\begin{Shaded}
\begin{Highlighting}[]
\NormalTok{Z\_\{N\_h\}(u)=\textbackslash{}frac\{1\}\{K\}\textbackslash{}sum\_\{m=0\}\^{}\{K{-}1\}e\^{}\{i(2m+1)u\}}
\end{Highlighting}
\end{Shaded}

The larger \texttt{N\_h} is, the narrower the localization width
becomes.

If the base wavelength of the parent system is \texttt{λ\_0} and the
required localization width is \texttt{Δx}, the required highest
odd-harmonic order is estimated as

\begin{Shaded}
\begin{Highlighting}[]
\NormalTok{N\_\{h,\textbackslash{}rm req\}\textbackslash{}sim\textbackslash{}frac\{\textbackslash{}lambda\_0\}\{\textbackslash{}Delta x\}}
\end{Highlighting}
\end{Shaded}

Similarly, for a temporal base period \texttt{T\_0} and temporal
localization width \texttt{Δt},

\begin{Shaded}
\begin{Highlighting}[]
\NormalTok{N\_\{h,\textbackslash{}rm req\}\^{}\{(t)\}\textbackslash{}sim\textbackslash{}frac\{T\_0\}\{\textbackslash{}Delta t\}}
\end{Highlighting}
\end{Shaded}

Thus, when the parent base wavelength or period is very large, extremely
large highest odd-harmonic order is required to construct local spatial
or temporal localization.

\subsection{Working Axiom C: Coupling
Entrance}\label{working-axiom-c-coupling-entrance}

Classification: Working hypothesis

The entrance to interaction is alignment of low-order base phases.

\subsection{Working Axiom D: Observation
Map}\label{working-axiom-d-observation-map}

Classification: Working hypothesis

\begin{Shaded}
\begin{Highlighting}[]
\NormalTok{s\_\textbackslash{}alpha=\textbackslash{}mathrm\{Loc\}\textbackslash{}left[\textbackslash{}int A(\textbackslash{}theta)M\_\textbackslash{}alpha(\textbackslash{}theta)d\textbackslash{}theta\textbackslash{}right]}
\end{Highlighting}
\end{Shaded}

\texttt{Loc} is a map that produces a stable localized output from a
correlation amplitude. Observation here is an internal definition of
this axiom system and does not claim connection to standard-theory
observation.

\subsection{Working Axiom E: Correlation-Strength
Squared}\label{working-axiom-e-correlation-strength-squared}

Classification: Working hypothesis

\begin{Shaded}
\begin{Highlighting}[]
\NormalTok{P(\textbackslash{}alpha)=\textbackslash{}frac\{|C\_\textbackslash{}alpha|\^{}2\}\{\textbackslash{}sum\_\textbackslash{}beta |C\_\textbackslash{}beta|\^{}2\}}
\end{Highlighting}
\end{Shaded}

This formula is a normalization of correlation readout strength within
this axiom system.

\subsection{Working Axiom F: Low-Order Base Synchronization
Sharing}\label{working-axiom-f-low-order-base-synchronization-sharing}

Classification: Working hypothesis

\begin{Shaded}
\begin{Highlighting}[]
\NormalTok{\textbackslash{}theta\_A{-}\textbackslash{}theta\_B=\textbackslash{}Delta}
\end{Highlighting}
\end{Shaded}

Low-order base synchronization sharing is a state in which two waveforms
share a low-order base phase difference.

\subsection{Working Axiom G: Loss of Synchronous
Demodulability}\label{working-axiom-g-loss-of-synchronous-demodulability}

Classification: Working hypothesis

\begin{Shaded}
\begin{Highlighting}[]
\NormalTok{V(t)=\textbackslash{}left|\textbackslash{}sum\_n w\_n e\^{}\{i\textbackslash{}epsilon\_n(t)\}\textbackslash{}right|}
\end{Highlighting}
\end{Shaded}

Phase-fluctuation average:

\begin{Shaded}
\begin{Highlighting}[]
\NormalTok{V(t)=\textbackslash{}left|\textbackslash{}sum\_n w\_n e\^{}\{{-}\textbackslash{}frac12\textbackslash{}sigma\_n\^{}2(t)\}\textbackslash{}right|}
\end{Highlighting}
\end{Shaded}

Loss of synchronous demodulability is a state in which phase
fluctuations prevent a correlation readout including higher orders from
producing a stable output.

\subsection{Working Axiom H: Hierarchy}\label{working-axiom-h-hierarchy}

Classification: Working hypothesis

A closed wave system has the hierarchy

\begin{Shaded}
\begin{Highlighting}[]
\NormalTok{low{-}order base mode}
\NormalTok{+ higher{-}order localized mode}
\end{Highlighting}
\end{Shaded}

\begin{center}\rule{0.5\linewidth}{0.5pt}\end{center}

\section{7. Readout, Memory, and Complex
Representation}\label{readout-memory-and-complex-representation}

\subsection{I/Q Readout}\label{iq-readout}

Classification: Theoretical insight

\begin{Shaded}
\begin{Highlighting}[]
\NormalTok{z=a+ib}
\end{Highlighting}
\end{Shaded}

\begin{Shaded}
\begin{Highlighting}[]
\NormalTok{a=r\textbackslash{}cos\textbackslash{}theta,\textbackslash{}qquad b=r\textbackslash{}sin\textbackslash{}theta}
\end{Highlighting}
\end{Shaded}

This is read as a pair of ninety-degree phase-shifted wave sources of
the same physical quantity.

\begin{Shaded}
\begin{Highlighting}[]
\NormalTok{u=r\textbackslash{}cos\textbackslash{}theta+r\textbackslash{}sin\textbackslash{}theta}
\NormalTok{=\textbackslash{}sqrt2 r\textbackslash{}cos\textbackslash{}left(\textbackslash{}theta{-}\textbackslash{}frac\{\textbackslash{}pi\}\{4\}\textbackslash{}right)}
\end{Highlighting}
\end{Shaded}

Complex representation is the minimal memory of a present value and a
ninety-degree delayed component.

\subsection{Harmonic Information
Readout}\label{harmonic-information-readout}

Classification: Theoretical insight

\begin{Shaded}
\begin{Highlighting}[]
\NormalTok{Z(\textbackslash{}theta)=\textbackslash{}sum\_n r\_n e\^{}\{i(n\textbackslash{}theta+\textbackslash{}phi\_n)\}}
\end{Highlighting}
\end{Shaded}

Readout map:

\begin{Shaded}
\begin{Highlighting}[]
\NormalTok{P(Z)=\textbackslash{}operatorname\{Re\}(Z)+\textbackslash{}operatorname\{Im\}(Z)}
\end{Highlighting}
\end{Shaded}

\begin{Shaded}
\begin{Highlighting}[]
\NormalTok{P(Z)=\textbackslash{}sum\_n \textbackslash{}sqrt2 r\_n}
\NormalTok{\textbackslash{}cos\textbackslash{}left(n\textbackslash{}theta+\textbackslash{}phi\_n{-}\textbackslash{}frac\{\textbackslash{}pi\}\{4\}\textbackslash{}right)}
\end{Highlighting}
\end{Shaded}

The order is preserved, and the readout phase is shifted by forty-five
degrees.

\subsection{Conservation Side and Readout
Side}\label{conservation-side-and-readout-side}

Classification: Theoretical insight

\begin{Shaded}
\begin{Highlighting}[]
\NormalTok{z\textbackslash{}bar z=a\^{}2+b\^{}2}
\end{Highlighting}
\end{Shaded}

is a conjugate norm.

\begin{Shaded}
\begin{Highlighting}[]
\NormalTok{z\^{}2=a\^{}2{-}b\^{}2+2iab}
\end{Highlighting}
\end{Shaded}

preserves interference phase.

Axiom 1 uses the \texttt{z\^{}2} type, not
\texttt{z\textbackslash{}bar\ z}.

\subsection{Square Register}\label{square-register}

Classification: Working hypothesis

Internal register:

\begin{Shaded}
\begin{Highlighting}[]
\NormalTok{\textbackslash{}mathcal\{A\}\_0=\textbackslash{}sum\_k A\_k\^{}2}
\end{Highlighting}
\end{Shaded}

External standard intensity:

\begin{Shaded}
\begin{Highlighting}[]
\NormalTok{a\_0=\textbackslash{}sum\_k |A\_k|\^{}2=\textbackslash{}sum\_k A\_k\textbackslash{}bar A\_k}
\end{Highlighting}
\end{Shaded}

\texttt{\textbackslash{}mathcal\{A\}\_0} preserves interference phase.
\texttt{a\_0} removes interference phase through conjugation.

\begin{center}\rule{0.5\linewidth}{0.5pt}\end{center}

\section{8. Statistical
Classification}\label{statistical-classification}

\subsection{Pair Phase Difference}\label{pair-phase-difference}

Classification: Definition

\begin{Shaded}
\begin{Highlighting}[]
\NormalTok{\textbackslash{}Delta\_\{ij\}=\textbackslash{}theta\_i{-}\textbackslash{}theta\_j}
\end{Highlighting}
\end{Shaded}

\begin{Shaded}
\begin{Highlighting}[]
\NormalTok{\textbackslash{}binom\{N\}\{2\}=\textbackslash{}frac\{N(N{-}1)\}\{2\}}
\end{Highlighting}
\end{Shaded}

Pair composition:

\begin{Shaded}
\begin{Highlighting}[]
\NormalTok{\textbackslash{}psi\_i+\textbackslash{}psi\_j}
\NormalTok{=2A\textbackslash{}cos\textbackslash{}frac\{\textbackslash{}Delta\_\{ij\}\}\{2\}e\^{}\{i(\textbackslash{}theta\_i+\textbackslash{}theta\_j)/2\}}
\end{Highlighting}
\end{Shaded}

Pair interference intensity:

\begin{Shaded}
\begin{Highlighting}[]
\NormalTok{I\_\{ij\}=4A\^{}2\textbackslash{}cos\^{}2\textbackslash{}frac\{\textbackslash{}Delta\_\{ij\}\}\{2\}}
\end{Highlighting}
\end{Shaded}

\subsection{Overlap Degree}\label{overlap-degree}

Classification: Definition

\begin{Shaded}
\begin{Highlighting}[]
\NormalTok{R\_N=\textbackslash{}left|\textbackslash{}frac\{A\_0\}\{N\}\textbackslash{}sum\_\{k=1\}\^{}\{N\}e\^{}\{i\textbackslash{}theta\_k\}\textbackslash{}right|}
\end{Highlighting}
\end{Shaded}

\begin{Shaded}
\begin{Highlighting}[]
\NormalTok{\textbackslash{}rho\_N=\textbackslash{}frac\{R\_N\}\{A\_0\}}
\NormalTok{=\textbackslash{}frac\{1\}\{N\}\textbackslash{}left|\textbackslash{}sum\_\{k=1\}\^{}\{N\}e\^{}\{i\textbackslash{}theta\_k\}\textbackslash{}right|}
\end{Highlighting}
\end{Shaded}

\begin{Shaded}
\begin{Highlighting}[]
\NormalTok{0\textbackslash{}le\textbackslash{}rho\_N\textbackslash{}le1}
\end{Highlighting}
\end{Shaded}

\subsection{Fermion Degree}\label{fermion-degree}

Classification: Definition

\begin{Shaded}
\begin{Highlighting}[]
\NormalTok{F\_N=}
\NormalTok{\textbackslash{}frac\{2\}\{N(N{-}1)\}}
\NormalTok{\textbackslash{}sum\_\{i\textless{}j\}\textbackslash{}sin\^{}2\textbackslash{}frac\{\textbackslash{}Delta\_\{ij\}\}\{2\}}
\end{Highlighting}
\end{Shaded}

\begin{Shaded}
\begin{Highlighting}[]
\NormalTok{B\_N=1{-}F\_N}
\end{Highlighting}
\end{Shaded}

\begin{Shaded}
\begin{Highlighting}[]
\NormalTok{B\_N=}
\NormalTok{\textbackslash{}frac\{2\}\{N(N{-}1)\}}
\NormalTok{\textbackslash{}sum\_\{i\textless{}j\}\textbackslash{}cos\^{}2\textbackslash{}frac\{\textbackslash{}Delta\_\{ij\}\}\{2\}}
\end{Highlighting}
\end{Shaded}

Relation to overlap degree:

\begin{Shaded}
\begin{Highlighting}[]
\NormalTok{F\_N=\textbackslash{}frac\{N\}\{2(N{-}1)\}(1{-}\textbackslash{}rho\_N\^{}2)}
\end{Highlighting}
\end{Shaded}

Therefore,

\begin{Shaded}
\begin{Highlighting}[]
\NormalTok{\textbackslash{}rho\_N=0\textbackslash{}Rightarrow F\_N=\textbackslash{}frac\{N\}\{2(N{-}1)\}}
\end{Highlighting}
\end{Shaded}

\texttt{B\_N} is the pair in-phase degree.

Classification:

\begin{Shaded}
\begin{Highlighting}[]
\NormalTok{F\_N=0: boson{-}like}
\NormalTok{N=2, F\_N=1: purely fermion{-}like}
\NormalTok{0\textless{}F\_N\textless{}1: elmion{-}like}
\end{Highlighting}
\end{Shaded}

\subsection{Boson-Like Composite Wave}\label{boson-like-composite-wave}

Classification: Definition

\begin{Shaded}
\begin{Highlighting}[]
\NormalTok{\textbackslash{}theta\_1=\textbackslash{}theta\_2=\textbackslash{}cdots=\textbackslash{}theta\_N}
\end{Highlighting}
\end{Shaded}

\begin{Shaded}
\begin{Highlighting}[]
\NormalTok{F\_N=0,\textbackslash{}qquad B\_N=1,\textbackslash{}qquad \textbackslash{}rho\_N=1,\textbackslash{}qquad R\_N=A\_0}
\end{Highlighting}
\end{Shaded}

\subsection{Purely Fermion-Like Composite
Wave}\label{purely-fermion-like-composite-wave}

Classification: Definition

\begin{Shaded}
\begin{Highlighting}[]
\NormalTok{N=2,\textbackslash{}qquad \textbackslash{}Delta\_\{12\}=\textbackslash{}pi}
\end{Highlighting}
\end{Shaded}

\begin{Shaded}
\begin{Highlighting}[]
\NormalTok{e\^{}\{i\textbackslash{}theta\_1\}+e\^{}\{i\textbackslash{}theta\_2\}=0}
\end{Highlighting}
\end{Shaded}

\begin{Shaded}
\begin{Highlighting}[]
\NormalTok{F\_2=1,\textbackslash{}qquad \textbackslash{}rho\_2=0,\textbackslash{}qquad R\_2=0}
\end{Highlighting}
\end{Shaded}

For three or more components, the condition that all pairs be opposite
in phase cannot hold.

\begin{Shaded}
\begin{Highlighting}[]
\NormalTok{\textbackslash{}Delta\_\{12\}=\textbackslash{}pi,\textbackslash{}quad \textbackslash{}Delta\_\{23\}=\textbackslash{}pi}
\NormalTok{\textbackslash{}Rightarrow}
\NormalTok{\textbackslash{}Delta\_\{13\}=2\textbackslash{}pi\textbackslash{}equiv0}
\end{Highlighting}
\end{Shaded}

\subsection{Elmion-Like Composite
Wave}\label{elmion-like-composite-wave}

Classification: Definition

The main condition is

\begin{Shaded}
\begin{Highlighting}[]
\NormalTok{0\textless{}F\_N\textless{}1}
\end{Highlighting}
\end{Shaded}

A subset with nonzero external composite amplitude is

\begin{Shaded}
\begin{Highlighting}[]
\NormalTok{0\textless{}\textbackslash{}rho\_N\textless{}1}
\end{Highlighting}
\end{Shaded}

\begin{Shaded}
\begin{Highlighting}[]
\NormalTok{\textbackslash{}rho\_N=0\textbackslash{}not\textbackslash{}Rightarrow F\_N=1}
\end{Highlighting}
\end{Shaded}

Zero external composite amplitude and pure two-body opposite-phase
character are not the same.

\end{document}
